\subsection{6.3  Lessons from the Case Study}

The case study illustrates several key properties of the APIP framework. First, the discipline of cause attribution prevents tier mismatch---applying the correct intervention tier is as important as applying any intervention. Second, the behavioral contract formalizes the trigger, making it objective and consistently applied rather than dependent on informal noticing. Third, the documented APIP record creates organizational memory: when Aria's RAG index next requires refreshing, the prior APIP provides a template for the intervention and a baseline for exit criteria. Fourth, the regression test requirement embedded in the intervention plan prevents the cascading failure that occurred in the ad-hoc scenario. Fifth, the Kirkpatrick-level evaluation at check-ins provides confidence that recovery is durable rather than superficial.

